\chapter{视界、状态方程与真空能}

\section{优选参考系、超光速退行与视界}

宇宙微波背景确实挑出了一个特殊的参考系。在那个使该辐射看起来各向同性的参考系中，附近的星系平均而言也处于静止。如果我们穿过背景运动，前向的光子会因多普勒效应而移向更高能量，后向的光子则移向更低能量。优选参考系就是使微波背景看起来各向同性的那个参考系，它同时也是平均而言本地星系处于静止的参考系。

\begin{classroomqa}
\lecturetimestamp{00:00:25}
\audiencequestion{宇宙辐射的表观波长是否依赖于观测者的运动，这是否意味着存在一个优选参考系？}
\lecturerresponse{是的。穿过微波背景的运动会产生多普勒各向异性，因此存在一个优选的宇宙学参考系：即微波背景看起来各向同性的那个参考系。它同时也是平均而言本地星系处于静止的参考系。}
\end{classroomqa}

第二个问题涉及退行速度。哈勃定律说，足够遥远的星系按照下式分离
\begin{equation}
v = H d .
\end{equation}
对于足够大的\(d\)，这一退行速度超过光速。这并不与相对论矛盾，因为该禁令是局部的：一个物体不能以超光速从就在我们眼前的另一个物体旁经过。关于遥远星系的陈述是关于相距遥远之处之间几何的陈述。

\begin{classroomqa}
\lecturetimestamp{00:04:48}
\audiencequestion{两个星系能否具有大于\(c\)的相对速度？}
\lecturerresponse{能，如果它们相距足够远。哈勃退行可以超过光速而不违反局部相对论，而退行速度达到\(c\)之处就是一个视界。}
\end{classroomqa}

一旦退行达到光速，朝我们发出的信号就不再取得净进展。该光线在局部仍以极限速度运动，但它必须克服的距离增长得太快。

\begin{classroomqa}
\lecturetimestamp{00:07:03}
\audiencequestion{朝相反方向发出的光线难道不仍会朝我们行进吗？}
\lecturerresponse{在局部它确实如此，但空间增长得比光线所能克服的更快。该光线相对于邻近点以极限速度运动，而回到我们这里的总距离却不断增加。}
\end{classroomqa}

\section{膨胀度规与直线类比}

一个有用的图像是一维的。想象一条由原子组成的直线。如果我们拉伸一条普通的材料直线，它最终会断裂。现在修改这个图像：每当相邻原子分离得太多时，我们就在它们之间再插入一个原子。直线不断增长，但它的局部性质并不会变得稀疏。

\begin{figure}
\centering
\includegraphics[width=0.60\textwidth]{lecture_04_figure_02.png}
\caption{用来描绘膨胀度规的一维直线。}
\end{figure}

在这个图像中仍然存在一个局部的极限信号速度，正如普通直线上的声音有一个极限速度一样。规则是局部的：相对于邻近点，扰动不能超过那个速度。然而，如果直线足够长，并且新点插入得足够快，那么它足够遥远的部分退行得如此之快，以至于信号永远无法返回。

因此，狭义相对论的速度相加只是故事的一部分。它适用于局部情形，即观测者在同一点彼此经过之时。对于膨胀几何中相距遥远的星系，它并不是合适的工具。

\begin{classroomqa}
\lecturetimestamp{00:09:42}
\audiencequestion{在这种情形下，通常的相对论速度相加定律是否仍然成立？}
\lecturerresponse{在局部，在一点处，是的。如果观测者一个接一个地彼此经过，通常的狭义相对论规则适用。但它并不是描述相距遥远、超过\(c\)的退行的正确方式；在那里必须回到膨胀空间的几何。}
\end{classroomqa}

同一图像给出了对红移的简单说明。在直线上画一个波。当直线膨胀时，波长随之膨胀。在膨胀几何中，这不仅仅是一个提示性的意象。膨胀时空中的波动方程确实会产生被拉伸的波长。

\begin{classroomqa}
\lecturetimestamp{00:14:42}
\audiencequestion{波随宇宙拉伸这一陈述是理论的、经验的，还是两者兼有？}
\lecturerresponse{它是膨胀时空中波传播的一个理论推论，同时也得到天文观测的支持。}
\end{classroomqa}

这个图像不是橡皮筋。橡皮筋被拉伸时会变细，而生活在这样一种介质中的人会看到局部性质发生变化。在这里，局部性质保持不变，而整体距离尺度却在增长。

\section{最后散射与可观测宇宙}

视界是退行速度达到光速的地方。最后散射面位于其稍内侧。从视界发出的信号会被无限拉伸；而最后散射辐射则具有巨大但有限的红移。

\begin{classroomqa}
\lecturetimestamp{00:17:53}
\audiencequestion{超光速退行、最后散射面与可观测宇宙之间是什么关系？}
\lecturerresponse{视界是退行速度达到光速的地方。最后散射面位于其稍内侧，因为它的红移虽然巨大但有限。就实际目的而言，它标志着可观测宇宙的边缘。}
\end{classroomqa}

最后散射辐射的红移大致为 \(10^3\) 倍。这已经如此之大，以至于在许多用途上，最后散射面都非常接近视界。可观测宇宙是宇宙中光能够到达我们的那部分。

在晚期，膨胀趋近于真空主导的 regime。此时哈勃量趋于常数，视界趋近于固定大小，星系和微波辐射都逐渐从视野中消失。

\begin{classroomqa}
\lecturetimestamp{00:22:43}
\audiencequestion{既然 \(H\) 随时间变化，视界距离是否也在变化，微波背景最终会消失吗？}
\lecturerresponse{在晚期 \(H\) 趋于常数，因此视界趋近于固定距离。微波背景持续被拉伸、冷却并逐渐消失。遥远的星系也会穿过视界并从视野中消失。}
\end{classroomqa}

只有已经受引力束缚的结构仍然可见。太阳系不会飞散，星系不会飞散，但大尺度宇宙会变成一个孤立的岛屿。

\section{基本膨胀方程与三种组分}

现在基本膨胀方程为
\begin{equation}
\left(\frac{\dot a}{a}\right)^2 = \frac{8\pi G}{3}\rho ,
\qquad c=1 .
\end{equation}
这里 \(a(t)\) 是标度因子，\(\rho\) 是总能量密度。

\begin{figure}
\centering
\includegraphics[width=0.74\textwidth]{lecture_04_figure_03.png}
\caption{基本膨胀方程以及一维线段类比。}
\end{figure}

有三种在动力学上重要的组分。

对于物质，
\begin{equation}
\rho=\frac{\text{const.}}{a^3},
\qquad
a(t)\sim t^{2/3}.
\end{equation}
这是有质量物体随整体膨胀运动、并且在共动参考系中基本静止的情形。

对于辐射，
\begin{equation}
\rho=\frac{\text{const.}}{a^4},
\qquad
a(t)\sim t^{1/2}.
\end{equation}
多出的 \(a\) 幂次来自宇宙膨胀时每个光子的红移。

还有
\begin{equation}
\rho_0=\text{constant},
\end{equation}
真空能，也称为暗能量，或者在不考虑约定数值因子时称为宇宙学常数。这一组分完全不会稀释。

下一步是把能量密度与压强联系起来。

\section{辐射压与冷物质}

从辐射开始。考虑一个一维盒子，也就是一段长度为 \(L\) 的线段，其中包含来回反弹的光子。在一维中，力与压强实际上是同一回事。

取一个动量大小为 \(p\) 的光子。当它从壁上反射时，其动量变化为
\begin{equation}
\Delta p = 2p.
\end{equation}
它与同一壁连续两次碰撞之间的时间为
\begin{equation}
\Delta t = \frac{2L}{c}.
\end{equation}
因此壁上的平均力为
\begin{equation}
F=\frac{\Delta p}{\Delta t}=\frac{pc}{L}.
\end{equation}
对于光子，
\begin{equation}
E=pc,
\end{equation}
所以
\begin{equation}
F=\frac{E}{L}.
\end{equation}
能量除以长度就是一维能量密度。在这个一维情形中，
\begin{equation}
P=\rho .
\end{equation}

对于许多光子，取平均后同样的关系仍然成立。总力就是单位长度的总能量。

现在转到三维。如果所有光子都只沿一个轴运动，它们就会只对一对表面施加压强，而对其他表面不施加压强。这是故意错误的，但它立刻显示出必须修正什么。在各向同性的光子浴中，动量分布在所有方向上。平均而言，能量密度的三分之一贡献给每个压强分量。结果是
\begin{equation}
P=\frac{\rho}{3}.
\end{equation}

对于冷物质，答案不同。取一个共动盒子，即随宇宙流动运动的一个小区域。其中的粒子携带能量，但如果它们相对于盒子静止，就不会撞到盒壁上。因此压强为
\begin{equation}
P=0.
\end{equation}

\begin{figure}
\centering
\includegraphics[width=0.74\textwidth]{lecture_04_figure_04.png}
\caption{物质与辐射：\(a\sim t^{2/3}\)、\(P=0\)，以及 \(a\sim t^{1/2}\)、\(P=\rho/3\)。}
\end{figure}

黑板上的总结是
\begin{align}
\text{matter:}\qquad & a(t)\sim t^{2/3}, \qquad P=0,\\
\text{radiation:}\qquad & a(t)\sim t^{1/2}, \qquad P=\frac{\rho}{3}.
\end{align}

\begin{classroomqa}
\lecturetimestamp{00:54:27}
\audiencequestion{在膨胀宇宙中，粒子实际上不是在彼此远离吗？}
\lecturerresponse{是的，但压强论证是针对一个无穷小共动盒子进行的。在这个局部尺度上，热力学推理仍然成立。}
\end{classroomqa}

\section{状态方程}

压强被认为与能量密度成正比，
\begin{equation}
P=w\rho .
\end{equation}
在黑板上，同一个参数写作\(\omega\)。对于两种基本情况，
\begin{equation}
w=0 \quad \text{for cold matter},
\qquad
w=\frac{1}{3} \quad \text{for radiation}.
\end{equation}

中间情况是可能的。介于\(0\)与\(1/3\)之间的值描述既不完全冷也不完全相对论性的物质。温暗物质就是这里提到的例子。

\begin{figure}
\centering
\includegraphics[width=0.74\textwidth]{lecture_04_figure_05.png}
\caption{状态方程\(P=\omega\rho\)连同密度标度论证。}
\end{figure}

下一个问题是，\(w\)是否控制\(\rho\)对\(a\)的依赖关系。答案来自一个热力学恒等式。

取一个体积为\(V\)的盒子。如果其内容物施加压强\(P\)，那么当盒子膨胀\(dV\)时，对壁所做的功为
\begin{equation}
dE=-P\,dV.
\end{equation}
现在把能量写成
\begin{equation}
E=\rho V.
\end{equation}
其微分为
\begin{equation}
dE=\rho\,dV+V\,d\rho .
\end{equation}
令两个表达式相等，得到
\begin{equation}
\rho\,dV+V\,d\rho=-P\,dV,
\end{equation}
或
\begin{equation}
V\,d\rho=-(\rho+P)\,dV.
\end{equation}
利用\(P=w\rho\)，
\begin{equation}
V\,d\rho=-(1+w)\rho\,dV.
\end{equation}
除以\(\rho V\)：
\begin{equation}
\frac{d\rho}{\rho}=-(1+w)\frac{dV}{V}.
\end{equation}
积分，
\begin{equation}
\rho \propto V^{-(1+w)}.
\end{equation}
对于共动宇宙学盒子，
\begin{equation}
V\propto a^3,
\end{equation}
因此
\begin{equation}
\rho \propto a^{-3(1+w)}.
\end{equation}

先前的情况立刻得到恢复：
\begin{align}
w=0 &\Longrightarrow \rho \propto a^{-3},\\
w=\frac{1}{3} &\Longrightarrow \rho \propto a^{-4}.
\end{align}

这就是\(w\)重要的原因。一旦\(\rho(a)\)已知，就可以把它代回弗里德曼方程。然后尝试幂律\(a\sim t^\alpha\)，并确定对于该状态方程宇宙如何膨胀。

\section{真空能与负压}

有一种特殊情况表现不同：
\begin{equation}
w=-1.
\end{equation}
于是
\begin{equation}
1+w=0,
\end{equation}
所以指数消失，并且
\begin{equation}
\rho=\rho_0=\text{constant}.
\end{equation}
相应的压强为
\begin{equation}
P=-\rho .
\end{equation}

负压本身并不神秘。正压向外推盒子的壁。负压向内拉。它是张力。拉伸的橡皮筋给出了日常例子：它储存正的能量，但它把两端拉在一起。

把常数密度\(\rho_0\)代入弗里德曼方程：
\begin{equation}
\left(\frac{\dot a}{a}\right)^2=\frac{8\pi G}{3}\rho_0.
\end{equation}
右边只是一个数，所以
\begin{equation}
\frac{\dot a}{a}=\sqrt{\frac{8\pi G}{3}\rho_0}.
\end{equation}
这就是哈勃量真正为常数的特殊情况。该方程可以写成
\begin{equation}
\dot a=\sqrt{\frac{8\pi G}{3}\rho_0}\,a,
\end{equation}
其解为
\begin{equation}
a(t)\propto \exp\!\left(\sqrt{\frac{8\pi G}{3}\rho_0}\,t\right).
\end{equation}

这是指数膨胀。相关的时间尺度是\(e\)折时间，而不是倍增时间，并且这里使用的粗略估计把它置于宇宙年龄的量级。除历史性的数值约定之外，真空能与宇宙学常数是同一物理成分。

\begin{classroomqa}
\lecturetimestamp{01:27:16}
\audiencequestion{如果真空能驱动指数膨胀，难道一切最终不会被撕裂吗？}
\lecturerresponse{不一定。束缚系统可以抵抗它。星系、太阳系或原子不必被撕开，而足够大且稀薄的结构确实参与膨胀并最终分离。}
\end{classroomqa}

\begin{classroomqa}
\lecturetimestamp{01:31:43}
\audiencequestion{当取弗里德曼方程的平方根时，负号中有什么有趣的东西吗？}
\lecturerresponse{有。平方根给出两个分支。一个是膨胀解，另一个是时间反演的收缩解。方程本身允许两者。}
\end{classroomqa}

\section{三个时代与观测到的膨胀历史}

三种成分都存在，并且它们以不同方式标度：
\begin{equation}
\begin{aligned}
\rho_{\text{radiation}}&\propto a^{-4},\\
\rho_{\text{matter}}&\propto a^{-3},\\
\rho_0&=\text{constant}.
\end{aligned}
\end{equation}
这立刻意味着三个时代。

在非常小的\(a\)下，辐射占主导。之后，物质占主导。在足够大的\(a\)下，真空能占主导。相应的膨胀定律为
\begin{equation}
a(t)\sim t^{1/2},
\qquad
a(t)\sim t^{2/3},
\qquad
a(t)\sim e^{Ht}.
\end{equation}

这里引用的粗略时间是：辐射在大约\(10^4\)年之后让位于物质。宇宙后来在大约\(10^5\)年之后变得透明。从物质主导到加速膨胀的转变发生得晚得多，在几十亿年之后，大约在所采用的量级估计中的宇宙历史中点。

这段历史是如何推断出来的？我们看不到未来，但能看到过去。看向更远的距离意味着看向更早的时间。通过观测遥远星系、测量它们的运动如何与距离相关，并重建哈勃量作为回溯时间的函数，人们重建了标度因子的时间历史。辐射主导时代太早，无法在普通星系天文学中直接看到，但物质主导时代以及向真空主导的转变在重建曲线中是可见的。

\section{遥远未来的最后散射}

未来的微波天空并不是星系通过视界消失的同一个故事。我们在更晚时间看到的辐射来自更远的物质，所以效应是看得更深、更冷，而不是看着同样的原子简单地退行离开。

\begin{classroomqa}
\lecturetimestamp{01:48:07}
\audiencequestion{如果最后散射面逐渐隐入视界，这与星系运动到视界之外是同一回事吗？}
\lecturerresponse{不是。后来发出的辐射并非来自我们现在所看到的辐射的同一批原子。它来自更远处的原子，所以这就像越来越深地看进宇宙，而辐射则显得越来越冷。}
\end{classroomqa}

\section{曲率、大挤压图景与晚期视界}

到此为止，曲率项一直被略去。将它恢复后得到
\begin{equation}
\left(\frac{\dot a}{a}\right)^2
=
\frac{8\pi G}{3}\rho
-
\frac{K}{a^2}.
\end{equation}
从数值上看，这一项似乎很小，但从概念上看，它很重要。在较老的牛顿语言中，它区分了总能量为正、为零和为负的情形。在相对论语言中，它区分了空间几何。

在真空能尚未被认真对待之前，人们常常设想一个物质主导的宇宙，其膨胀最终会反转。其结果就是旧的大挤压图景：从大爆炸开始膨胀，然后重新收缩。其数学仍然有效，但它似乎不再描述我们的宇宙。

\begin{classroomqa}
\lecturetimestamp{01:53:51}
\audiencequestion{这种真空能与量子真空能有关吗？}
\lecturerresponse{是的，或者至少人们相信如此。}
\end{classroomqa}

今天的粗略组成为
\begin{equation}
\frac{\rho_0}{\rho_{\text{total}}}\sim 0.7,
\qquad
\frac{\rho_{\text{matter}}}{\rho_{\text{total}}}\sim 0.3.
\end{equation}
在该物质成分中，发光物质约占总量的\(0.05\)，暗物质约占\(0.25\)。辐射在今天可以忽略不计。暗物质与发光物质以相同方式标度，因此即使物质、辐射和真空的相对丰度随时代而变化，它们的比例仍保持固定。

在真空主导的极限下，
\begin{equation}
H=\sqrt{\frac{8\pi G}{3}\rho_0}.
\end{equation}
于是哈勃定律变为
\begin{equation}
v=Hd.
\end{equation}
视界由\(v=1\)定义，因此其距离为
\begin{equation}
d_H=\frac{1}{H}
=
\frac{1}{\sqrt{8\pi G\rho_0/3}}.
\end{equation}
这里给出的粗略数值约为\(10\)到\(12\)亿光年。

重要之处在于晚期行为。一旦膨胀变为指数式，视界距离就变得固定。星系继续穿过它，并因红移越来越大而消失。微波背景也同样红移远去。在遥远的未来，只有附近受引力束缚的结构仍然可见，宇宙看起来就像一座孤岛。

\section{从最后散射面看几何}

最后，几何可以从最后散射面来检验。假设到该表面的距离已知。再假设能够在那里辨认出一个其物理大小已知得足够好的团块。那么就可以将该大小和距离与观测到的角大小进行比较。

在平坦空间中，角度是欧几里得角度。在正曲率空间中，它会显得更大；在负曲率空间中，则更小。类比就是普通的测量：取一个已知距离处已知大小的物体，并将测得的角度与欧几里得几何进行比较。

\begin{classroomqa}
\lecturetimestamp{02:08:09}
\audiencequestion{那里有什么物体的大小是我们已知的例子？}
\lecturerresponse{微波背景中最大的团块。我们认为，从它们演化所需的时间以及声速可以知道它们的大小。}
\end{classroomqa}

这里相关的对象是微波背景上最大的团块。它们观测到的角大小与欧几里得几何相符，误差在百分之几以内。据此，空间看起来是平坦的。

\section{测得的\(w\)与最后一个问题}

\(w\)在观测上的重要性现在清楚了。如果天文学家报告\(w\)非常接近\(-1\)，他们就是在说真空能成分几乎恒定。这正是这里所强调的要点。

这并不意味着广义相对论迫使\(w\)必须恰好等于\(-1\)。这里的要点是经验性的。讲座简要指出，原则上并不排除更激进的可能性，但有充分理由相信的是真空能几乎恒定的情形。

\begin{classroomqa}
\lecturetimestamp{02:11:04}
\audiencequestion{在物质和辐射的密度计算中，我们是否假设了能量是恒定的，然后调整标度因子来得出公式？如果前两种情况确实如此，那么在真空能情形中会发生什么？在那里密度保持恒定。}
\lecturerresponse{不。我并没有假设它是恒定的。我们用的是能量变化为\(-P\,dV\)。对于一个随膨胀一起运动的假想共动盒子，在壁上做了功，这就解释了能量的变化。真空能情形以不同的方式提出了这个问题，而这里的结束语是：宇宙学中的能量守恒语言并没有说明太多，因为宇宙的总能量为零。}
\end{classroomqa}